\documentclass[11pt]{article}

\usepackage[margin=1in]{geometry}
\usepackage{amsmath,amssymb}
\usepackage{graphicx}
\usepackage{booktabs}
\usepackage{hyperref}
\usepackage[dvipsnames]{xcolor}
\usepackage{minted}

% Minted settings
\setminted{fontsize=\footnotesize, breaklines, frame=single, framesep=2mm}
\setminted[yaml]{style=friendly}
\setminted[python]{style=friendly}
\setminted[julia]{style=friendly}

\title{Supplementary Materials:\\LLM-Assisted Extraction of QSP Calibration Targets}
\author{Joel Eliason}
\date{}

\begin{document}

\maketitle

\tableofcontents
\newpage

%=============================================================================
\section{S1: Complete SubmodelTarget Example}
\label{sec:complete-example}
%=============================================================================

This section presents a complete calibration target for activated pancreatic stellate cell (aPSC) proliferation, demonstrating all schema components.

\subsection{Overview}

This target constrains the proliferation rate parameter \texttt{k\_apsc\_prolif} using data from Schneider et al. (2001), who measured PSC proliferation in response to PDGF stimulation via BrdU incorporation assay.

\subsection{Complete YAML Specification}

\begin{minted}{yaml}
# Schema: SubmodelTarget
# Target: Activated PSC proliferation rate

target_id: psc_proliferation_PDAC_deriv001

# =============================================================================
# INPUTS: Extracted values with full provenance
# =============================================================================
inputs:
  - name: pdgf_fold_mean
    value: 4.37
    units: dimensionless
    input_type: direct_measurement
    role: target
    source_ref: Schneider2001_PSC_PDGF
    source_location: Abstract
    value_snippet: "Cell proliferation (4.37 +/- 0.49- and 2.96 +/-
                    0.39-fold of control)."

  - name: pdgf_fold_sd
    value: 0.49
    units: dimensionless
    input_type: direct_measurement
    role: auxiliary
    source_ref: Schneider2001_PSC_PDGF
    source_location: Abstract
    value_snippet: "Cell proliferation (4.37 +/- 0.49- and 2.96 +/-
                    0.39-fold of control)."

# =============================================================================
# CALIBRATION: Everything needed for inference
# =============================================================================
calibration:
  parameters:
    - name: k_apsc_prolif
      units: 1/day
      prior:
        distribution: lognormal
        mu: 0.0      # log(1.0) - expect ~1/day for 4x fold-change in 1 day
        sigma: 1.0
        rationale: |
          Wide log-normal prior centered at 1/day. A 4-fold increase in
          1 day implies k = ln(4) = 1.4/day, so prior median of 1/day
          is reasonable.

  state_variables:
    - name: aPSC_rel
      units: dimensionless
      initial_condition:
        value: 1.0
        rationale: |
          PSC proliferation data are reported as fold of control. We
          normalize the initial aPSC population to 1.0 (100% control)
          to match this convention.

  model:
    type: exponential_growth
    rate_constant: k_apsc_prolif
    data_rationale: |
      BrdU incorporation assay measures proliferation as fold-change
      over ~24 hours under constant PDGF stimulus. With a single high
      PDGF dose and short duration, exponential growth (dN/dt = k*N)
      fits the data structure.
    submodel_rationale: |
      The full QSPIO_PDAC model has aPSC proliferation via:
      k_apsc_prolif * aPSC * (PDGF / (PDGF_50_prolif + PDGF)).
      At 50 ng/mL PDGF >> PDGF_50_prolif, the saturation term ~ 1,
      reducing to exponential growth with rate k_apsc_prolif.

  independent_variable:
    name: time
    units: day
    span: [0.0, 1.0]
    rationale: |
      Proliferation assay duration assumed to be ~24 hours based on
      typical cytokine incubation protocols for rat PSCs.

  measurements:
    - name: proliferation_fold_change
      observable:
        type: identity
        state_variables: [aPSC_rel]
        rationale: |
          The state variable aPSC_rel is the fold-change in aPSC number
          relative to baseline (control = 1.0), which directly matches
          the reported proliferation fold-change.
      units: dimensionless
      uses_inputs:
        - pdgf_fold_mean
        - pdgf_fold_sd
      evaluation_points: [1.0]
      sample_size: 3
      sample_size_rationale: |
        Sample size not explicitly reported; assumed n=3 independent
        experiments, typical for in vitro PSC BrdU proliferation assays.
      distribution_code: |
        def derive_distribution(inputs, ureg):
            import numpy as np

            rng = np.random.default_rng(42)

            mean = inputs['pdgf_fold_mean'].magnitude
            sd = inputs['pdgf_fold_sd'].magnitude
            n = 10000

            # Use lognormal for positive fold-changes
            mu_log = np.log(mean**2 / np.sqrt(mean**2 + sd**2))
            sigma_log = np.sqrt(np.log(1 + sd**2 / mean**2))

            samples = rng.lognormal(mu_log, sigma_log, n)

            median_val = np.median(samples)
            ci_lower = np.percentile(samples, 2.5)
            ci_upper = np.percentile(samples, 97.5)

            return {
                'median': [median_val],
                'ci95': [[ci_lower, ci_upper]],
                'units': 'dimensionless'
            }
      likelihood:
        distribution: lognormal
        rationale: |
          Fold-changes are positive and multiplicative, making
          lognormal appropriate.

  identifiability_notes: |
    With a single timepoint (~24h) and normalized initial condition
    (aPSC_rel(0)=1), only k_apsc_prolif is estimated. If substantial
    PSC death occurred over the assay interval, the inferred rate
    would represent net growth rather than pure proliferation.

# =============================================================================
# EXPERIMENTAL CONTEXT
# =============================================================================
experimental_context:
  species: rat
  system: in_vitro_primary_cells
  indication: PDAC
  cell_types:
    - name: pancreatic stellate cell
      phenotype: activated (alpha-SMA positive)
      isolation_method: Isolated from normal Wistar rat pancreas,
                        activated by culture on plastic
  culture_conditions:
    culture_type: 2d_monolayer

# =============================================================================
# STUDY INTERPRETATION
# =============================================================================
study_interpretation: |
  Schneider et al. (2001) measured proliferation of cultured rat
  pancreatic stellate cells (PSCs) in response to growth factors using
  BrdU incorporation. Addition of 50 ng/mL PDGF increased PSC
  proliferation to 4.37 +/- 0.49-fold of control, indicating that PDGF
  acts as a strong mitogen for activated PSCs. We interpret the
  50 ng/mL PDGF condition as approximating a near-maximal PDGF stimulus
  and use the reported fold-change over ~24 hours as a quantitative
  constraint on the maximal PDGF-driven proliferation rate parameter
  k_apsc_prolif.

key_assumptions:
  - |
    50 ng/mL PDGF corresponds to a near-saturating concentration for
    PSC proliferation, so PDGF/(PDGF_50_prolif + PDGF) ~ 1.
  - |
    The proliferation assay duration is approximately 24 hours,
    consistent with typical cytokine incubation protocols.
  - |
    BrdU incorporation fold-change is proportional to fold-change in
    PSC cell number (negligible death over 24 hours).
  - |
    Rat activated PSC proliferation kinetics in vitro provide an upper
    bound for human tumor aPSC proliferation in vivo.

key_study_limitations:
  - |
    Proliferation was measured via BrdU incorporation rather than
    direct cell counts.
  - |
    Single PDGF concentration and single assay interval constrain
    only the maximal rate.
  - |
    Data from rat PSCs in vitro, not human PDAC-associated stellate cells.
  - |
    Background serum proliferative signals not explicitly modeled.

# =============================================================================
# DATA SOURCES
# =============================================================================
primary_data_source:
  doi: "10.1152/ajpcell.2001.281.2.C532"
  title: "Identification of mediators stimulating proliferation and
          matrix synthesis of rat pancreatic stellate cells"
  authors: [Schneider, Schmid, Liptay, Schmid, Pohl]
  year: 2001
  source_tag: Schneider2001_PSC_PDGF

secondary_data_sources:
  - doi: "10.1136/gut.44.4.534"
    title: "Pancreatic stellate cells are activated by proinflammatory
            cytokines: implications for pancreatic fibrogenesis"
    authors: [Apte, Haber, Darby, Rodgers, McCaughan, Korsten, Pirola, Wilson]
    year: 1999
    source_tag: Apte1999_PSC_Cytokines
    contribution: |
      Provides context on PSC activation and cytokine response protocols.
\end{minted}

\subsection{Schema Structure Summary}

The SubmodelTarget schema has the following top-level structure:

\begin{itemize}
    \item \textbf{target\_id}: Unique identifier for the calibration target
    \item \textbf{inputs}: List of values extracted from literature with provenance
    \item \textbf{calibration}: Everything needed for inference code generation
    \begin{itemize}
        \item \texttt{parameters}: Parameters to estimate with priors
        \item \texttt{state\_variables}: ODE state variables with initial conditions
        \item \texttt{model}: Model type specification (built-in or custom)
        \item \texttt{independent\_variable}: Time or other independent variable
        \item \texttt{measurements}: Observables with likelihoods
        \item \texttt{identifiability\_notes}: Discussion of parameter identifiability
    \end{itemize}
    \item \textbf{experimental\_context}: Species, system, cell types, culture conditions
    \item \textbf{study\_interpretation}: Scientific interpretation of the data
    \item \textbf{key\_assumptions}: Assumptions made in using this data
    \item \textbf{key\_study\_limitations}: Known limitations of the source study
    \item \textbf{primary\_data\_source}: Primary literature reference with DOI
    \item \textbf{secondary\_data\_sources}: Additional references
\end{itemize}

%=============================================================================
\section{S2: Supported Model Types}
\label{sec:model-types}
%=============================================================================

The schema supports the following built-in model types:

\begin{table}[htbp]
\centering
\caption{Built-in model types with their ODE formulations.}
\label{tab:model-types}
\begin{tabular}{lll}
\toprule
Model Type & ODE & Parameters \\
\midrule
\texttt{exponential\_growth} & $dy/dt = k \cdot y$ & \texttt{rate\_constant} \\
\texttt{first\_order\_decay} & $dy/dt = -k \cdot y$ & \texttt{rate\_constant} \\
\texttt{logistic} & $dy/dt = k \cdot y \cdot (1 - y/K)$ & \texttt{rate\_constant}, \texttt{carrying\_capacity} \\
\texttt{saturation} & $dy/dt = k \cdot (1 - y)$ & \texttt{rate\_constant} \\
\texttt{two\_state} & $dA/dt = -k \cdot A$, $dB/dt = k \cdot A$ & \texttt{forward\_rate} \\
\texttt{michaelis\_menten} & $dy/dt = -V_{max} \cdot y / (K_m + y)$ & \texttt{vmax}, \texttt{km} \\
\texttt{direct\_conversion} & (analytical formula) & \texttt{formula} \\
\texttt{direct\_fit} & (curve fitting) & \texttt{curve} \\
\texttt{custom} & (user-provided) & \texttt{code}, \texttt{code\_julia} \\
\bottomrule
\end{tabular}
\end{table}

%=============================================================================
\section{S3: Validation Checks}
\label{sec:validators}
%=============================================================================

The validation framework includes ten automated checks implemented as Pydantic \texttt{@model\_validator} methods. Each validator runs after field parsing and raises a \texttt{ValidationError} with detailed messages when constraints are violated.

\subsection{Summary of Validators}

\begin{enumerate}
    \item \textbf{Input reference validation}: All \texttt{uses\_inputs} references match defined input names
    \item \textbf{Source reference validation}: All \texttt{source\_ref} values match defined source tags
    \item \textbf{Parameter reference validation}: Model parameter roles reference defined parameters
    \item \textbf{State variable reference validation}: Observable \texttt{state\_variables} match defined state variables
    \item \textbf{DOI resolution}: All DOIs resolve via CrossRef API
    \item \textbf{Value-in-snippet validation}: Extracted values appear verbatim in quoted snippets
    \item \textbf{Unit validation}: All unit strings parse with Pint
    \item \textbf{Code syntax validation}: Custom code passes Python AST parsing
    \item \textbf{Span ordering validation}: Time spans have start $<$ end
    \item \textbf{ODE model requirements}: ODE models have required state variables and spans
\end{enumerate}

\subsection{DOI Resolution Validator}

The DOI validator queries the CrossRef API to verify that every cited DOI resolves to a valid publication. This catches fabricated citations before they enter the calibration pipeline. The validator also performs fuzzy matching on title and year to detect metadata inconsistencies.

\begin{minted}[fontsize=\scriptsize]{python}
def resolve_doi(doi: str) -> Optional[dict]:
    """Resolve DOI and get metadata from CrossRef."""
    doi_clean = doi.replace("https://doi.org/", "").strip()

    url = f"https://doi.org/{doi_clean}"
    headers = {"Accept": "application/vnd.citationstyles.csl+json"}
    response = requests.get(url, headers=headers, timeout=10)

    if response.status_code != 200:
        return None  # DOI does not resolve

    metadata = response.json()
    return {
        "title": metadata.get("title", ""),
        "first_author": metadata.get("author", [{}])[0].get("family", ""),
        "year": metadata.get("issued", {}).get("date-parts", [[]])[0][0]
    }

@model_validator(mode="after")
def validate_doi_resolution_and_metadata(self) -> "SubmodelTarget":
    """Validate that DOIs resolve and metadata matches."""
    errors = []

    metadata = resolve_doi(self.primary_data_source.doi)
    if metadata is None:
        errors.append(f"DOI '{self.primary_data_source.doi}' failed to resolve")
    else:
        # Check title match (fuzzy)
        if not fuzzy_match(self.primary_data_source.title, metadata["title"]):
            errors.append(f"Title mismatch with CrossRef")
        # Check year match
        if self.primary_data_source.year != metadata["year"]:
            errors.append(f"Year mismatch: {self.primary_data_source.year} vs {metadata['year']}")

    if errors:
        raise ValueError("DOI validation errors:\n  - " + "\n  - ".join(errors))
    return self
\end{minted}

\textbf{Example failure}: If an LLM fabricates a citation with DOI \texttt{10.1234/fake.2023.12345}, the CrossRef query returns a 404 error and validation fails with:
\begin{verbatim}
ValidationError: DOI '10.1234/fake.2023.12345' failed to resolve.
Verify at https://doi.org/10.1234/fake.2023.12345
\end{verbatim}

\subsection{Value-in-Snippet Validator}

The value-in-snippet validator checks that each extracted numeric value appears in its associated source text. This catches hallucinations where an LLM extracts a plausible but incorrect value. The validator handles multiple numeric formats including scientific notation, percentages, and Unicode superscripts.

\begin{minted}[fontsize=\scriptsize]{python}
def check_value_in_text(text: str, value: float) -> bool:
    """Check if numeric value appears in text."""
    text_norm = text.lower().replace(",", "")
    patterns = []

    # Direct value and common formats
    patterns.append(str(value))
    patterns.append(f"{value:.1f}")
    patterns.append(f"{value:.2f}")

    # Scientific notation (e.g., "1.5e-3", "1.5E-03")
    if abs(value) < 0.01 or abs(value) > 1000:
        patterns.append(f"{value:e}")
        patterns.append(f"{value:.2e}")

    # Percentage format (if value is between 0 and 1)
    if 0 < value < 1:
        pct = value * 100
        patterns.extend([f"{pct}%", f"{pct:.0f}%", f"{pct:.1f}%"])

    return any(p.lower() in text_norm for p in patterns)

@model_validator(mode="after")
def validate_input_values_in_snippets(self) -> "SubmodelTarget":
    """Validate that extracted values appear in their value_snippet."""
    errors = []
    for inp in self.inputs:
        if not inp.value_snippet:
            continue
        # Skip types that may not have explicit values in text
        if inp.input_type in (InputType.INFERRED_ESTIMATE, InputType.ASSUMED_VALUE):
            continue

        if not check_value_in_text(inp.value_snippet, inp.value):
            errors.append(
                f"Input '{inp.name}': value {inp.value} not found in snippet "
                f"'{inp.value_snippet[:60]}...'"
            )

    if errors:
        raise ValueError("Value-snippet mismatches (possible hallucination):\n  - "
                        + "\n  - ".join(errors))
    return self
\end{minted}

\textbf{Example failure}: If an input claims \texttt{value: 4.37} but the snippet reads ``Cell proliferation was 3.21-fold of control'', validation fails with:
\begin{verbatim}
ValidationError: Value-snippet mismatches (possible hallucination):
  - Input 'pdgf_fold_mean': value 4.37 not found in snippet
    'Cell proliferation was 3.21-fold of control...'
\end{verbatim}

This validator is particularly effective because LLMs that hallucinate values rarely fabricate consistent snippets---the mismatch provides a reliable detection signal.

\subsection{Unit Validation}

The unit validator parses all unit strings using Pint, a Python library for physical quantities. Malformed or unrecognized units trigger validation failure.

\begin{minted}[fontsize=\scriptsize]{python}
@model_validator(mode="after")
def validate_units_are_valid_pint(self) -> "SubmodelTarget":
    """Validate that all unit strings are valid Pint units."""
    from pint import UnitRegistry
    ureg = UnitRegistry()
    errors = []

    def check_unit(unit_str: str, location: str):
        try:
            ureg(unit_str)
        except Exception as e:
            errors.append(f"{location}: '{unit_str}' is not valid ({e})")

    for inp in self.inputs:
        check_unit(inp.units, f"Input '{inp.name}'")
    for param in self.calibration.parameters:
        check_unit(param.units, f"Parameter '{param.name}'")

    if errors:
        raise ValueError("Invalid Pint units:\n  - " + "\n  - ".join(errors))
    return self
\end{minted}

\textbf{Example failure}: If a parameter specifies \texttt{units: cells/mL/day} (invalid compound unit), validation fails with:
\begin{verbatim}
ValidationError: Invalid Pint units:
  - Parameter 'k_recruit': 'cells/mL/day' is not valid (...)
\end{verbatim}

\subsection{Reference Consistency Validators}

Four validators check that all cross-references within the schema are consistent:

\begin{itemize}
    \item \textbf{Input references}: Every \texttt{uses\_inputs} entry and \texttt{initial\_condition.input\_ref} must match a defined input name.
    \item \textbf{Source references}: Every \texttt{source\_ref} must match \texttt{primary\_data\_source.source\_tag} or a \texttt{secondary\_data\_sources[].source\_tag}.
    \item \textbf{Parameter references}: When a model field like \texttt{rate\_constant} contains a string (not an \texttt{InputRef}), it must match a parameter name in \texttt{calibration.parameters}.
    \item \textbf{State variable references}: Every state variable referenced in \texttt{observable.state\_variables} must exist in \texttt{calibration.state\_variables}.
\end{itemize}

These validators catch internal inconsistencies where an LLM generates references to entities that don't exist elsewhere in the document.

%=============================================================================
\section{S4: Generated Julia Code Structure}
\label{sec:julia-code}
%=============================================================================

The translator generates Julia scripts with the following structure:

\begin{minted}{julia}
using DifferentialEquations
using Turing
using Distributions

# Observed data (from inputs layer)
const obs_proliferation_median = 4.37
const obs_proliferation_sigma = 0.25  # derived from CI95

# ODE function
function ode_proliferation!(du, u, p, t)
    k = p[1]
    du[1] = k * u[1]  # exponential growth
end

# Simulate function
function simulate_proliferation(k; t_span=(0.0, 1.0), u0=[1.0])
    prob = ODEProblem(ode_proliferation!, u0, t_span, [k])
    sol = solve(prob, Tsit5())
    return sol[end][1]  # final value
end

# Turing model
@model function proliferation_model()
    # Prior
    k_apsc_prolif ~ LogNormal(0.0, 1.0)

    # Simulate
    pred = simulate_proliferation(k_apsc_prolif)

    # Likelihood
    obs_proliferation_median ~ LogNormal(log(pred), obs_proliferation_sigma)
end

# Run inference
model = proliferation_model()
chain = sample(model, NUTS(), MCMCThreads(), 1000, 4)
\end{minted}

For joint inference over multiple targets, the translator identifies shared parameters by name and generates a single model with one prior per unique parameter.

%=============================================================================
\section{S5: Pydantic Model Definitions}
\label{sec:pydantic-models}
%=============================================================================

This section provides the key Pydantic model definitions that implement the SubmodelTarget schema. The full implementation is available in the repository at \texttt{src/qsp\_llm\_workflows/core/calibration/submodel\_target.py}.

\subsection{Top-Level Model}

The \texttt{SubmodelTarget} class is the entry point for validation. When a YAML file is loaded, Pydantic recursively validates all nested models.

\begin{minted}[fontsize=\scriptsize]{python}
class SubmodelTarget(BaseModel):
    """
    Submodel-based calibration target for QSP parameter inference.

    Separates:
    - `inputs`: What was extracted from papers (with full provenance)
    - `calibration`: How to use those inputs for inference
    """
    target_id: str = Field(description="Unique identifier for this target")

    # Extracted data with provenance
    inputs: List[Input] = Field(
        description="Values extracted from literature with full provenance"
    )

    # Calibration specification
    calibration: Calibration = Field(
        description="Everything needed for inference code generation"
    )

    # Experimental context
    experimental_context: ExperimentalContext = Field(
        description="Experimental context for the target"
    )

    # Study interpretation
    study_interpretation: str = Field(
        description="Scientific interpretation of how the study data informs the model"
    )
    key_assumptions: List[str] = Field(
        description="Key assumptions made in using this data"
    )
    key_study_limitations: Optional[List[str]] = Field(
        default=None, description="Known limitations of the study"
    )

    # Data sources
    primary_data_source: PrimaryDataSource = Field(
        description="Primary literature source"
    )
    secondary_data_sources: Optional[List[SecondaryDataSource]] = Field(
        default=None, description="Additional literature sources"
    )

    # Validators (10 total) defined via Pydantic decorators
    # See Section S3 for implementation details
\end{minted}

\subsection{Input Model}

Each input captures a single value extracted from literature with full provenance:

\begin{minted}[fontsize=\scriptsize]{python}
class InputType(str, Enum):
    """Type of input extracted from literature."""
    DIRECT_MEASUREMENT = "direct_measurement"   # Value reported directly
    PROXY_MEASUREMENT = "proxy_measurement"     # Requires conversion
    EXPERIMENTAL_CONDITION = "experimental_condition"  # Protocol choice
    INFERRED_ESTIMATE = "inferred_estimate"     # Interpreted from qualitative text
    ASSUMED_VALUE = "assumed_value"             # Domain knowledge, not in source

class InputRole(str, Enum):
    """Role of input in calibration."""
    INITIAL_CONDITION = "initial_condition"  # IC for ODE integration
    TARGET = "target"                        # Calibration target (likelihood term)
    FIXED_PARAMETER = "fixed_parameter"      # Fixed value in model
    AUXILIARY = "auxiliary"                  # Supporting data

class Input(BaseModel):
    """A value extracted from literature with full provenance."""
    name: str = Field(description="Unique identifier for this input")
    value: float = Field(description="Extracted numeric value")
    units: str = Field(description="Units of the value")
    uncertainty: Optional[Uncertainty] = Field(default=None)
    n: Optional[int] = Field(default=None, description="Sample size")
    input_type: InputType
    role: Optional[InputRole] = Field(default=None)
    source_ref: str = Field(description="Reference to source_tag in data sources")
    source_location: str = Field(description="Location within the source")
    value_snippet: Optional[str] = Field(
        default=None,
        description="Exact text from paper containing the value (for validation)"
    )
\end{minted}

\subsection{Model Type Discriminated Union}

The schema uses Pydantic's discriminated union pattern to support multiple model types, each with type-specific required fields:

\begin{minted}[fontsize=\scriptsize]{python}
class BaseModelSpec(BaseModel):
    """Base class for all model specifications."""
    data_rationale: str = Field(
        description="Why this model type fits the experimental data"
    )
    submodel_rationale: str = Field(
        description="Why this is a valid submodel of the full QSP model"
    )

class ExponentialGrowthModel(BaseModelSpec):
    """Exponential growth: dy/dt = k * y"""
    type: Literal["exponential_growth"] = "exponential_growth"
    rate_constant: ParameterRole = Field(
        description="Rate constant parameter name or input_ref"
    )

class TwoStateModel(BaseModelSpec):
    """Two-state transition: A -> B with first-order kinetics."""
    type: Literal["two_state"] = "two_state"
    forward_rate: ParameterRole = Field(
        description="Forward transition rate constant"
    )

class CustomModel(BaseModelSpec):
    """Custom ODE with user-provided code"""
    type: Literal["custom"] = "custom"
    code: str = Field(description="Python ODE function")
    code_julia: str = Field(description="Julia ODE function for inference")

# Discriminated union selects model class based on 'type' field
Model = Annotated[
    Union[
        FirstOrderDecayModel,
        ExponentialGrowthModel,
        LogisticModel,
        MichaelisMentenModel,
        TwoStateModel,
        SaturationModel,
        DirectConversionModel,
        DirectFitModel,
        CustomModel,
    ],
    Field(discriminator="type"),
]
\end{minted}

The \texttt{discriminator="type"} annotation tells Pydantic to inspect the \texttt{type} field in the YAML to determine which model class to instantiate. This ensures that an \texttt{exponential\_growth} model is validated against \texttt{ExponentialGrowthModel} (which requires \texttt{rate\_constant}) while a \texttt{two\_state} model is validated against \texttt{TwoStateModel} (which requires \texttt{forward\_rate}).

\subsection{Prior Specification}

Priors are specified with distribution type and parameters:

\begin{minted}[fontsize=\scriptsize]{python}
class PriorDistribution(str, Enum):
    """Supported prior distribution types."""
    LOGNORMAL = "lognormal"    # For positive parameters (rates, densities)
    NORMAL = "normal"          # For unconstrained parameters
    UNIFORM = "uniform"        # For bounded parameters
    HALF_NORMAL = "half_normal"  # For positive parameters with mode at 0

class Prior(BaseModel):
    """Prior distribution specification for Bayesian inference."""
    distribution: PriorDistribution
    mu: Optional[float] = Field(default=None, description="Location parameter")
    sigma: Optional[float] = Field(default=None, description="Scale parameter")
    lower: Optional[float] = Field(default=None, description="Lower bound")
    upper: Optional[float] = Field(default=None, description="Upper bound")
    rationale: Optional[str] = Field(default=None)

    @model_validator(mode="after")
    def validate_prior_params(self) -> "Prior":
        """Validate that required parameters are provided for each distribution."""
        if self.distribution == PriorDistribution.LOGNORMAL:
            if self.mu is None or self.sigma is None:
                raise ValueError("lognormal prior requires mu and sigma")
        elif self.distribution == PriorDistribution.UNIFORM:
            if self.lower is None or self.upper is None:
                raise ValueError("uniform prior requires lower and upper")
            if self.lower >= self.upper:
                raise ValueError("uniform lower must be < upper")
        return self
\end{minted}

\subsection{Loading and Validating YAML Files}

To load and validate a calibration target:

\begin{minted}[fontsize=\scriptsize]{python}
import yaml
from qsp_llm_workflows.core.calibration import SubmodelTarget

# Load YAML file
with open("psc_proliferation_PDAC_deriv001.yaml") as f:
    data = yaml.safe_load(f)

# Parse and validate (raises ValidationError on failure)
target = SubmodelTarget(**data)

# Access validated fields
print(f"Target: {target.target_id}")
print(f"Parameters: {[p.name for p in target.calibration.parameters]}")
print(f"Model type: {target.calibration.model.type}")
\end{minted}

If validation fails, Pydantic raises a \texttt{ValidationError} with detailed error messages:

\begin{verbatim}
pydantic_core._pydantic_core.ValidationError: 2 validation errors for SubmodelTarget
calibration.measurements.0.uses_inputs.0
  Input 'typo_input_name' not found in inputs [...]
primary_data_source.doi
  DOI '10.1234/fake' failed to resolve [...]
\end{verbatim}

\end{document}
